\chapter{State of the Art}
\label{ch:sota}

Whereas the fundamentals establish the terminology of safety assurance, safety
argumentation, and safety performance indicators, this chapter reviews the
\textit{current} state of research on their combined use for operational safety
assurance of \ac{ADS}. The review is organised from the general to the specific. It
begins with structured safety argumentation and the shift towards operational
assurance, narrows to \ac{SPI} based safety performance monitoring, then examines the
accident causation and investigation methods that supply the reasoning for deriving
indicators, surveys how mature safety domains transform lagging indicators into
leading ones, reviews the surrogate safety metrics that serve as candidate signals,
and closes with a consolidated statement of the research gap that motivates the
methodology in Chapter~\ref{ch:methodology}. Each section concludes with the aspect
that remains open.

\section{Structured Safety Argumentation for Automated Driving Systems}
\label{sec:sota_argumentation}

Safety assurance cases and the Goal Structuring Notation are introduced in
Section~\ref{sec:fund_gsn}; this section reviews how they are treated in current
research on automated driving. Among the available notations, namely the
Claims-Arguments-Evidence notation, the Structured Assurance Case Metamodel, and the
\ac{GSN}, the \ac{GSN} is the most widely used in safety-critical domains, and its
decomposition of a top-level goal into sub-goals provides the claim hierarchy to
which indicators are attached in this work \cite{kelly2004gsn}.

The assurance case for an \ac{ADS} is valid at the time of its submission to a
certification authority, based on the evidence available during development
\cite{laxman2024}. For automated driving, however, new or contradictory information
can emerge after deployment, revealing emergent system properties, environmental
uncertainties, or divergences in component performance that undermine the validity
of the original argument \cite{laxman2025}. Assurance cases are intended to be living
documents, yet in conventional practice they are treated as static artefacts that
are signed off once \cite{laxman2025}. The central limitation of pure design-time
argumentation is therefore that it cannot, on its own, account for the open and
dynamic operating environment of an \ac{ADS}.

\section{Operational Safety Assurance}
\label{sec:sota_operational}

Operational safety assurance addresses the limitation above by maintaining confidence
in the validity of the assurance case after deployment, using the growing
availability of field data \cite{laxman2024}. Research in this direction is gaining
traction, particularly for runtime model-based assurance cases. Wei et al. argue for
the transition from design-time to runtime assurance-case processes on the basis of
the Structured Assurance Case Metamodel \cite{wei2024access}, and Asaadi et al.
introduce dynamic assurance cases as a pathway to trusted autonomy for systems based
on artificial intelligence and machine learning \cite{asaadi2020dynamic}. Carlan et
al. extend this line of work to a dynamic safety-case management system for frontier
artificial-intelligence systems \cite{carlan2024dynamic}. These approaches establish
that the assurance case is to be treated as a dynamic artefact rather than a fixed
document.

A second line of research improves the rigour of the argument itself so that its
runtime challenges become explicit. Assurance~2.0 retains the structure of
traditional assurance cases but emphasises rigorous evidence assessment and the
active search for defeaters in order to counter confirmation bias
\cite{bloomfield2020assurance2}. Eliminative argumentation provides the underlying
mechanism, in which claims are recursively challenged by doubts, termed defeaters,
and confidence in a claim increases as these defeaters are resolved by additional
argumentation or evidence \cite{goodenough2015ea}. Three types of defeater are
distinguished, namely rebutting, undermining, and undercutting defeaters
\cite{laxman2025}. Diemert et al. build on this notion with incremental assurance, in
which confidence in an eliminative argument is raised by supplying additional
evidence that resolves the recorded defeaters \cite{diemert2023incremental}. Hawkins
and Ryan Conmy complement design-time argumentation with dialectic analysis, which
identifies how an argument may be challenged at runtime and derives corresponding
runtime monitoring requirements in a systematic manner \cite{hawkins2023dialectic}.
The common contribution of these approaches is a structured record of doubt. The open
aspect is that the resolution of many defeaters requires operational data that is
unavailable before deployment.

The integration of field data into this process is the specific concern of the
SafetyDevOps paradigm. The SafeOps approach applies the principles of automation,
feature-driven development, and operational monitoring to the iterative extension of
safety-critical products under ISO~26262 \cite{munk2022safeops}. Building on this,
Laxman et al. propose a cumulative operational safety assurance approach in which risk
assessment, safety measures, requirements, and their validation are built up
iteratively over time together with field data, so that the assurance case is
enhanced continuously rather than fixed at release \cite{laxman2024}. The authors
identify the absence of systematic methods for leveraging runtime vehicle data in a
cumulative manner as a principal research gap \cite{laxman2024}. This position
establishes the process context of the present work, but it does not by itself
specify which quantities are to be monitored.

\section{Safety Performance Indicators}
\label{sec:sota_spi}

Safety performance indicators, defined in Section~\ref{sec:fund_spi}, are the
mechanism by which field data is connected to the claims of an assurance case, and
their use is mandated for automated driving by ANSI/UL~4600 in its Section~16
\cite{ul4600} and by BSI PAS~1881 \cite{pas1881}, with complementary field-monitoring
obligations arising from ISO~26262 \cite{iso26262} and ISO~21448 \cite{iso21448}. The
research question is therefore not whether indicators are required, which the standards
settle, but how a useful set of them is obtained.

The distinction between leading and lagging indicators (Section~\ref{sec:fund_spi}) is
the organising property of the research. Diaz and Woon observe that the lower the level
of the safety-case claim to which an indicator is attached, the more likely it is to be
a leading indicator, because low-level claims are violated earlier than the high-level
outcomes they support \cite{diaz2022drspi}. This relationship between claim depth and
indicator lead time is illustrated in Figure~\ref{fig:spi_spectrum} and is central to
the derivation strategy pursued in this work. ANSI/UL~4600 reflects the same priority,
stating in its guidance that reliance on lagging indicators alone is not acceptable,
yet the standard provides only examples of indicators and leaves their derivation to
the practitioner \cite{ul4600}.

% ---------------------------------------------------------------------------
\begin{figure}[htbp]
    \centering
    \includegraphics[width=15cm]{spi_leading_lagging_spectrum}
    \caption{Relationship between the depth of a safety-case claim and the leading or
    lagging character of the attached safety performance indicator. Indicators on deep
    sub-claims tend to be leading and provide earlier warning, whereas indicators on
    high-level claims tend to be lagging (based on \cite{diaz2022drspi} and
    \cite{koopman2022howsafe}).}
    \label{fig:spi_spectrum}
\end{figure}
% ---------------------------------------------------------------------------

Several frameworks build on this leading-indicator property. The Dynamic Realtime
\ac{SPI} framework of Diaz and Woon attaches indicators computed at runtime to
sufficiently low safety-case sub-claims, so that a threshold violation signals that
the \ac{ADS} is encountering an untested situation, and assigns each such indicator an
improvement mechanism that adjusts vehicle behaviour to re-establish the validity of
the sub-claim \cite{diaz2022drspi}. Ratiu et al. address a complementary concern,
namely that the rationale by which a chosen set of \acp{SPI} contributes to safety is
often left implicit. They propose an explicit argument pattern for the use of
\acp{SPI}, analyse the defeaters that can undermine confidence in a chosen indicator
set, and introduce meta-indicators that monitor whether the indicators themselves
continue to perform as intended \cite{ratiu2024}. The systematic derivation of
indicators from an assurance case is addressed most directly by Laxman et al., whose
four-step approach represents the uncertainty of assurance-case claims using
Dempster-Shafer theory, propagates this uncertainty to the top goal, resolves claims
through eliminative argumentation, and defines \acp{SPI} for those claims that cannot
be resolved without runtime data \cite{laxman2025}. To the knowledge reported in that
work, no standard or otherwise widely used systematic approach for specifying
\acp{SPI} exists \cite{laxman2025}. The open aspect is that the derivation of a
\emph{leading} indicator from a \emph{lagging} one is not itself formalised, and the
reasoning that justifies why a given signal predicts a given loss event is drawn from the
accident investigation literature reviewed next.

\section{Accident Causation Models and Investigation Methods}
\label{sec:sota_causation}

The derivation of a leading indicator from a loss event rests on a model of how that
loss event is caused, and such models originate in accident investigation, which is
inherently a reactive activity performed after a loss. Three families of causation
model are relevant. The first is the sequential, component-failure family, of which
Fault Tree Analysis is the archetype. Fault Tree Analysis is a deductive, top-down
method that starts from an undesired top event and uses Boolean logic gates,
principally AND and OR, to decompose it into intermediate events and finally basic
events with associated failure data \cite{ericson1999,nasaftahandbook}. It is
codified in IEC~61025 \cite{iec61025} and, for automotive systems, recommended by
ISO~26262 for the higher automotive safety integrity levels \cite{iso26262}. The
top-down AND and OR decomposition supplies the structural skeleton for deriving
indicators, because a loss event treated as a top event decomposes into the
subsystem-level causes and observable signals that precede it. The recognised
limitation of classical Fault Tree Analysis is its assumption of discrete,
independent component failures, which does not capture the functional insufficiencies
of a machine-learning perception stack, where a component operates as specified yet
the specification itself is inadequate \cite{ruijters2015}. The structure is shown
schematically in Figure~\ref{fig:fault_tree}: Boolean reduction of the tree yields
minimal cut sets, the smallest combinations of basic events sufficient to cause the
top event, and these map directly onto the candidate signal sets from which leading
indicators are drawn.

% ---------------------------------------------------------------------------
\begin{figure}[htbp]
    \centering
    \includegraphics[width=13cm]{fault_tree_example}
    \caption{Deductive top-down structure of a fault tree: an undesired top event is
    decomposed through OR and AND gates into intermediate events and basic events
    (recreated, based on the fault-tree method of IEC~61025 \cite{iec61025}).}
    \label{fig:fault_tree}
\end{figure}
% ---------------------------------------------------------------------------

The second family is the epidemiological, barrier-based model. Reason's organisational
accident model distinguishes active failures from latent conditions and holds that an
accident occurs when weaknesses in successive defensive layers momentarily align
\cite{reason1990,reason2000}. The bowtie method operationalises this view by combining
a fault tree of threats leading to a central loss event with an event tree of
consequences, and by making the preventive and mitigating barriers explicit
\cite{ccps2018bowties}. The health of each barrier can then be monitored by a
dedicated leading indicator, which is the conceptual bridge from a causation model to
an \ac{SPI}. Application of the bowtie method to road-traffic and cooperative-driving
safety exists but remains uncommon \cite{ehlers2017}, and a comprehensive treatment
for a full \ac{ADS} safety case is not established. The relationship between the two
views is shown in Figure~\ref{fig:bowtie}, where the bowtie places the fault tree and
the event tree on either side of a central loss event, with the barriers whose health
a leading indicator monitors sitting on the preventive side.

% ---------------------------------------------------------------------------
\begin{figure}[htbp]
    \centering
    \includegraphics[width=15cm]{swiss_cheese_bowtie}
    \caption{Barrier-based accident model. Left: the Swiss Cheese view, in which a loss
    event occurs when weaknesses in successive defensive layers align. Right: the
    bowtie, combining a fault tree of threats and an event tree of consequences around
    a central loss event, with preventive and mitigating barriers made explicit
    (recreated, based on \cite{reason2000} and \cite{ccps2018bowties}).}
    \label{fig:bowtie}
\end{figure}
% ---------------------------------------------------------------------------

The third family is the systemic, control-theoretic model. The Systems-Theoretic
Accident Model and Processes reframes safety as the enforcement of constraints within
a hierarchical control structure rather than as a chain of component failures
\cite{leveson2011}. Its hazard-analysis technique, System-Theoretic Process Analysis,
identifies unsafe control actions and their causal scenarios \cite{stpahandbook}, and
its investigation counterpart, Causal Analysis based on System Theory, reconstructs
the control-structure failures behind an accident \cite{casthandbook}. Unsafe control
actions are candidate leading indicators, and the approach captures the non-failure,
functional-insufficiency hazards that Fault Tree Analysis misses under the safety of
the intended functionality \cite{iso21448}. The open aspect is that System-Theoretic
Process Analysis and Causal Analysis based on System Theory are qualitative and do
not natively produce thresholds or continuously monitored indicators
\cite{khastgir2021}. The hierarchical control view is shown in
Figure~\ref{fig:stamp}, in which each control action and feedback path is a place where
an unsafe control action, and hence a candidate leading indicator, can be located.

% ---------------------------------------------------------------------------
\begin{figure}[htbp]
    \centering
    \includegraphics[width=11cm]{stamp_control_structure}
    \caption{Hierarchical control structure of the Systems-Theoretic Accident Model and
    Processes. Safety is treated as the enforcement of constraints through control
    actions, with feedback informing each controller; inadequate control actions are
    the candidate sites for leading indicators (recreated, based on \cite{leveson2011}).}
    \label{fig:stamp}
\end{figure}
% ---------------------------------------------------------------------------

A distinct contribution comes from formal incident analysis. Why-Because Analysis
constructs a directed acyclic causal graph in which a factor is admitted as a cause
only if it passes a counterfactual test, that is, whether the effect would still have
occurred had the factor been absent \cite{ladkin1998}. This counterfactual test is
functionally identical to the prevention test used in the methodology of this work to
identify a proximate cause, and it therefore provides a formal, citable grounding for
that step. The open aspect common to all of these methods is that they are
retrospective and manual, applied to a single incident after the fact, and none has
been operationalised to generate forward-looking, continuously monitored indicators
for an \ac{ADS}.

\section{From Reactive Investigation to Proactive Monitoring}
\label{sec:sota_reactive_proactive}

The methods of Section~\ref{sec:sota_causation} are reactive, since they analyse a loss
that has already occurred, and the value of an \ac{SPI} lies in converting their
insight into a proactive signal that acts before a loss. Precursor analysis supplies the
theory for this conversion. It treats near-misses and off-nominal events as truncated
accident sequences that carry predictive information about rare, high-consequence
events, and uses Bayesian methods to update accident-likelihood estimates from
abundant low-severity data \cite{bieryi1995,saleh2013}. A leading indicator is, in
these terms, a formalised precursor, and precursor-estimation theory offers a
principled route to calibrating indicator thresholds from sparse loss data. The
associated caution is equally important. The fixed-ratio safety pyramid of Heinrich,
which assumes a stable numerical relationship between minor and major events
\cite{heinrich1931}, is contested by later evidence. Yorio and Moore, analysing 27,446
mining establishments over the period 2000 to 2012, find no significant difference in
the probability of a subsequent fatal event between establishments with few and with
zero lost or restricted days \cite{yoriomoore2018}. The quantitative link between a
precursor and a loss event must therefore be validated empirically rather than assumed.

Aviation demonstrates the reactive-to-proactive transition in mature form. Its safety
management system institutionalises a continuous loop that monitors indicators,
detects deviations, manages risk, and verifies the result \cite{icao9859}, and Flight
Operational Quality Assurance mines routine flight data for exceedances that act as
precursors before an accident occurs. This is the same lagging-to-leading
transformation that the present work systematises for automated driving, and it is
widely regarded as the origin of the \ac{SPI} concept later adopted by ANSI/UL~4600.
Within automated driving, the equivalent sentinel signal is the activation of
Autonomous Emergency Braking, a last-resort intervention triggered when normal control
has failed to preserve an adequate margin, typically parameterised by time-to-collision
\cite{eNCAP_aeb}. A rising activation rate is an early-warning precursor of collision
risk. The published treatment of Autonomous Emergency Braking is, however, oriented
towards component testing under rating protocols rather than towards a monitored fleet
indicator with a calibrated threshold feeding a safety case.

The lagging indicators of automated driving are documented in crash and disengagement
data. Analyses of the California disengagement and collision reports report, for
example, an average of one accident per 178 disengagements up to July 2017
\cite{favaro2018}, and the reporting mandated by the United States National Highway
Traffic Safety Administration provides a growing but incomplete record of automation-
involved crashes \cite{nhtsa_sgo}. The investigation of the Tesla Autopilot design by
that agency reviewed 956 crashes in which the system was alleged to have been engaged,
analysed 467 after exclusions, identified at least 13 involving one or more
fatalities, and concluded that the design contributed to foreseeable misuse and
avoidable crashes \cite{nhtsa_ea22002}. The same investigation records that only about
18 per cent of police-reported crashes involve an airbag deployment and that the
manufacturer largely receives data only for such crashes, which exposes the sparsity
and reporting bias of field lagging data \cite{nhtsa_ea22002}. This limitation is the
direct justification for generating leading indicators in simulation, as pursued in
this work, rather than relying on sparse field data.

\section{Deriving Leading Indicators from Lagging Indicators}
\label{sec:sota_derivation}

Mature safety domains beyond automated driving have converged on the practice of
pairing leading and lagging indicators, which frames the specific gap this work
addresses. In process safety, the guidance of the Center for Chemical Process Safety
and the recommended practice API~RP~754 define a tiered structure in which the higher
tiers are lagging outcome events and the lower tiers are leading indicators of barrier
and system performance, derived from the weaknesses of the barriers in a Swiss Cheese
or bowtie model \cite{ccps2018metrics,apirp754}. The step-by-step guide HSG254 of the
United Kingdom Health and Safety Executive adds the concept of dual assurance, pairing
a leading and a lagging indicator for each critical risk-control system
\cite{hse2006hsg254}. Across these domains the distinction is consistent, with lagging
indicators being retrospective, few, and severe, and leading indicators being
forward-looking and numerous.

Table~\ref{tab:crossdomain} sets the four domains side by side, which makes the common
structure and the common omission visible in one view.

% ---------------------------------------------------------------------------
\begin{table}[htbp]
    \centering
    \caption{Pairing of leading and lagging indicators across mature safety domains,
    showing the source of the leading indicator in each domain and whether a
    reproducible derivation procedure is published.}
    \label{tab:crossdomain}
    \small
    \begin{tabular}{@{}p{2.7cm} p{3.1cm} p{4.0cm} p{3.4cm}@{}}
        \toprule
        \textbf{Domain} & \textbf{Lagging indicator} & \textbf{Source of the leading indicator} & \textbf{Derivation procedure} \\
        \midrule
        Process safety \cite{ccps2018metrics,apirp754} & tiered loss-of-containment
        events & barrier weaknesses in a bowtie model & tiered structure given;
        derivation by expert judgement \\[2pt]
        Occupational safety \cite{hse2006hsg254} & lost-time and reportable
        incidents & risk-control systems, paired by dual assurance & pairing
        prescribed; derivation by expert judgement \\[2pt]
        Aviation \cite{icao9859} & accidents and serious incidents & routine
        flight-data exceedances & continuous monitoring institutionalised;
        derivation by expert judgement \\[2pt]
        Automated driving \cite{ul4600,ratiu2024} & collisions and
        disengagements & safety-case claims and surrogate metrics & none published;
        stated as an open gap \\
        \bottomrule
    \end{tabular}
\end{table}
% ---------------------------------------------------------------------------

The consistent shortfall is that the derivation itself is driven by expert judgement
rather than by a reproducible procedure. In every case a practitioner identifies
barriers from a causation model and then defines indicators of barrier health, and no
domain publishes a validated, step-by-step method that takes a specific lagging
indicator and produces calibrated leading indicators with explicit observability and
independence checks. For automated driving the gap is stated directly in the
literature. Ratiu et al. observe that there is no systematic discussion of the
challenges of deploying \acp{SPI} and frame their own contribution as only an initial
step towards a systematic framework \cite{ratiu2024}, and Diaz and Woon classify each
indicator as leading or lagging case by case rather than deriving them systematically
\cite{diaz2022drspi}. This absence of a reproducible derivation procedure is the
precise gap that the methodology of this work fills.

\section{Safety Metrics as Candidate Indicators}
\label{sec:sota_metrics}

The metric component of an \ac{SPI} is drawn from the body of surrogate safety
metrics defined in Section~\ref{sec:fund_metrics}, and the suitability of a metric
determines the quality of the resulting indicator. Fraade-Blanar et al. provide a
framework for measuring automated-vehicle safety that organises these measures along a
leading-to-lagging axis and catalogues their respective strengths and weaknesses
\cite{fraadeblanar2018}, and Westhofen et al. survey the wider landscape of criticality
metrics and analyse their suitability for different conflict types \cite{westhofen2023}.
The value of these metrics as established, pre-validated loss predictors means that the
derivation of this work can treat the signals as known and claim novelty in the
derivation procedure rather than in the predictive power of each signal.

\section{Conflicts in the Target Manoeuvre}
\label{sec:sota_conflicts}

The manoeuvre evaluated in this work, an unprotected left turn, folds together three
conflicts of different character, and the literature treats each with different metrics
and reports different limitations for each. Reviewing them separately establishes why a
single derivation method must be tested against all three rather than assumed to
transfer.

The first is the oncoming-vehicle conflict, a longitudinal left-turn-across-path
interaction. Studies of automated vehicles at unsignalised intersections report that
this conflict is among the most demanding, because the decision window is short and the
turn crosses the path of oncoming traffic \cite{leftturn_unsignalized}. Field testing
of driver-assistance systems in naturalistic left-turn-across-path conditions reports
that automatic emergency braking failed to activate braking early enough even when the
detection criteria were technically met, with decision windows on the order of three
seconds \cite{leftturn_field2025}. This indicates that a purely detection-based
indicator offers poor lead-time adequacy in this conflict.

The second is the pedestrian-crossing conflict, a crossing interaction for which
post-encroachment time rather than pure time-to-collision is the natural surrogate
measure \cite{allen1978pet}. Pedestrian protection is assessed in rating protocols
through dedicated automatic-emergency-braking pedestrian tests \cite{eNCAP_aeb}, yet
studies of such systems at intersections show that occlusion, in which the pedestrian
is hidden until late, sharply reduces the available response time and the achievable
avoidance rate \cite{ped_aeb_occluded}. The adequacy of a crossing indicator therefore
depends strongly on when the crossing road user first becomes observable, which differs
in character from the oncoming conflict.

The third is the rear-end, or following, conflict, in which the ego vehicle must
maintain a safe separation from a vehicle behind it while it decelerates for the turn.
The established surrogate measures here are time headway and the Responsibility-Sensitive
Safety safe distance \cite{shalev2017rss}, and recent work estimates rear-end collision
risk for automated vehicles in mixed traffic from real-world perception data
\cite{rearend_av2025}. The controlling quantity is the follower's ability to respond to
the ego vehicle's deceleration, which is again a different mechanism from the other two
conflicts.

Taken together, the three conflicts have distinct geometries and distinct
controlling quantities, so the adequacy of a given metric as a leading indicator is
conflict-dependent. This is the direct motivation for deriving and evaluating indicators
separately by conflict actor in Chapter~\ref{ch:methodology} and
Chapter~\ref{ch:results}, rather than assuming one indicator serves the whole manoeuvre.

\section{Research Gap}
\label{sec:sota_gap}

The reviewed literature supports a single, focused research gap. Structured
argumentation and operational-assurance approaches
(Section~\ref{sec:sota_argumentation}, Section~\ref{sec:sota_operational}) establish
that an assurance case must be maintained with field data through a SafetyDevOps
process. The \ac{SPI} literature (Section~\ref{sec:sota_spi}) establishes the
metric-threshold definition of an indicator, the value of leading indicators on deep
sub-claims, and a first systematic method for identifying which claims require
indicators. The accident-causation and cross-domain literature
(Section~\ref{sec:sota_causation} to Section~\ref{sec:sota_derivation}) supplies the
reasoning components, namely fault-tree decomposition, the counterfactual prevention
test, precursor theory, and barrier-health monitoring, but no domain integrates them
into a reproducible procedure. The metric literature (Section~\ref{sec:sota_metrics})
provides candidate signals and evidence that their adequacy depends on the conflict
type.

Three deficiencies remain. First, no reproducible method derives a set of leading
\acp{SPI} from a lagging \ac{SPI} for an \ac{ADS}; the derivation is left to expert
judgement across every domain, and the gap is stated explicitly for automated driving
\cite{ratiu2024,ul4600}. Second, indicators derived from an assurance case are, to the
extent reported, not evaluated quantitatively against simulation data, so their
ability to detect a developing violation early enough to be actionable, as distinct
from merely detecting that a violation has occurred, is not demonstrated. Third, the
sparsity and reporting bias of field lagging data \cite{favaro2018,nhtsa_ea22002} mean
that the quantitative link between a precursor and a loss event, which the Heinrich critique
shows must not be assumed \cite{yoriomoore2018}, has not been validated for the
relevant conflicts.

This thesis addresses these deficiencies with a systematic causal-chain derivation
method that takes a lagging \ac{SPI}, identifies its proximate cause through a
counterfactual prevention test, decomposes it with fault-tree gate logic to the
subsystem level, extracts the earliest observable signals, checks their observability
and independence, calibrates signal-level and rate-level thresholds, and attaches the
resulting leading \acp{SPI} to a \ac{GSN} safety case modelled in a \ac{GSN} modelling tool.
The indicators are evaluated from simulation data generated in the CARLA simulator
\cite{dosovitskiy2017carla} with an explicit distinction between detection and
lead-time adequacy, separated by conflict actor. The continuous assurance loop within
which the derived indicators operate is summarised in
Figure~\ref{fig:assurance_loop} and frames the methodology that follows.

% ---------------------------------------------------------------------------
\begin{figure}[htbp]
    \centering
    \includegraphics[width=15cm]{spi_assurance_loop}
    \caption{Continuous operational safety assurance loop. Safety performance
    indicators derived from the claims of the assurance case are monitored against
    field or simulation data; a threshold violation triggers a change-impact analysis
    and an update of the assurance case within a SafetyDevOps process (based on
    \cite{laxman2024} and \cite{laxman2025}).}
    \label{fig:assurance_loop}
\end{figure}
% ---------------------------------------------------------------------------